% !TEX program = xelatex
% Lernplatt - by Can Siebert
% Kurs 22 (Bo 143) - Tag 8 - Lehrskript
% Dynamische Preisstrategien & integrierte Plattform-Verkaufsarchitektur
%
% Self-contained xelatex source. Kein biblatex, keine externen Bilder.

\documentclass[11pt,a4paper,oneside]{article}

% --- Encoding & Sprache --------------------------------------------------
\usepackage{polyglossia}
\setdefaultlanguage{german}

% --- Geometrie -----------------------------------------------------------
\usepackage[a4paper,top=2cm,bottom=2cm,outer=2.5cm,inner=2.5cm,bindingoffset=1.5cm]{geometry}

% --- Fonts ---------------------------------------------------------------
\usepackage{fontspec}
\defaultfontfeatures{Ligatures=TeX}

\IfFontExistsTF{Charter}{%
  \setmainfont{Charter}[%
    Numbers=OldStyle,%
    UprightFeatures   = {SmallCapsFeatures={Letters=SmallCaps}}%
  ]%
}{%
  \IfFontExistsTF{Bitstream Charter}{%
    \setmainfont{Bitstream Charter}%
  }{%
    \setmainfont{TeX Gyre Schola}%
  }%
}

\IfFontExistsTF{Source Sans 3}{%
  \setsansfont{Source Sans 3}%
}{%
  \IfFontExistsTF{Source Sans Pro}{%
    \setsansfont{Source Sans Pro}%
  }{%
    \setsansfont{TeX Gyre Heros}%
  }%
}

\newcommand{\caveatfontfallback}{\itshape}
\IfFontExistsTF{Caveat}{%
  \newfontfamily\caveatfont{Caveat}[Scale=1.15]%
}{%
  \newcommand\caveatfont{\caveatfontfallback}%
}

\setmonofont{Latin Modern Mono}

% --- Mikro-Typografie ----------------------------------------------------
\usepackage{microtype}
\linespread{1.15}

% --- Farben & Boxen ------------------------------------------------------
\usepackage{xcolor}
\definecolor{lpInk}{RGB}{20,28,38}
\definecolor{lpAccent}{RGB}{22,90,140}
\definecolor{lpAccentLight}{RGB}{210,228,240}
\definecolor{lpAmber}{RGB}{222,138,30}
\definecolor{lpAmberLight}{RGB}{255,243,217}
\definecolor{lpAmberBorder}{RGB}{196,118,18}
\definecolor{lpGray}{RGB}{96,104,114}
\definecolor{lpGrayLight}{RGB}{238,240,243}

\usepackage[most]{tcolorbox}
\tcbuselibrary{breakable,skins}

\newtcolorbox{eselsbox}[1][Eselsbrücke]{%
  enhanced, breakable,
  colback=lpAmberLight,
  colframe=lpAmberBorder,
  boxrule=0.6pt,
  arc=2pt,
  borderline={0.4pt}{0pt}{lpAmberBorder, dashed},
  left=10pt,right=10pt,top=8pt,bottom=8pt,
  fonttitle=\sffamily\bfseries\color{lpAmberBorder},
  coltitle=lpAmberBorder,
  title={#1},
  attach boxed title to top left={xshift=8pt,yshift=-8pt},
  boxed title style={size=small, colback=lpAmberLight, colframe=lpAmberBorder, boxrule=0.4pt, arc=1pt},
  top=14pt
}

\newtcolorbox{merkbox}[1][Merksatz]{%
  enhanced, breakable,
  colback=lpAccentLight,
  colframe=lpAccent,
  boxrule=0.6pt,
  arc=2pt,
  left=10pt,right=10pt,top=8pt,bottom=8pt,
  fonttitle=\sffamily\bfseries\color{white},
  coltitle=white,
  title={#1},
  attach boxed title to top left={xshift=8pt,yshift=-8pt},
  boxed title style={size=small, colback=lpAccent, colframe=lpAccent, boxrule=0.4pt, arc=1pt},
  top=14pt
}

% --- Listen --------------------------------------------------------------
\usepackage{enumitem}
\setlist{itemsep=2pt, topsep=4pt, parsep=0pt}
\setlist[itemize,1]{label={\textbullet},leftmargin=1.6em}
\setlist[itemize,2]{label={\textendash},leftmargin=1.4em}
\setlist[enumerate,1]{label=\arabic*., leftmargin=1.8em}

% --- Tabellen ------------------------------------------------------------
\usepackage{tabularx}
\usepackage{array}
\usepackage{booktabs}
\usepackage{longtable}

% --- Kopf-/Fusszeilen ----------------------------------------------------
\usepackage{fancyhdr}
\usepackage{lastpage}
\pagestyle{fancy}
\fancyhf{}
\renewcommand{\headrulewidth}{0.3pt}
\renewcommand{\footrulewidth}{0pt}
\fancyhead[L]{\footnotesize\sffamily\color{lpGray}Lernplatt \textperiodcentered{} Kurs 22 \textperiodcentered{} Tag 8}
\fancyhead[R]{\footnotesize\sffamily\color{lpGray}Dynamische Preise \& Performance}
\fancyfoot[C]{\footnotesize\sffamily\color{lpGray}\thepage{} / \pageref{LastPage}}

\fancypagestyle{plain}{%
  \fancyhf{}%
  \renewcommand{\headrulewidth}{0pt}%
  \fancyfoot[C]{\footnotesize\sffamily\color{lpGray}\thepage{} / \pageref{LastPage}}%
}

% --- Section-Styling -----------------------------------------------------
\usepackage[explicit]{titlesec}

\titleformat{\section}[block]
  {\sffamily\Large\bfseries\color{lpAccent}}
  {\thesection.}{0.6em}{#1}
  [{\vspace{2pt}\color{lpAccent}\titlerule[0.6pt]}]

\titleformat{\subsection}[block]
  {\sffamily\large\bfseries\color{lpInk}}
  {\thesubsection}{0.6em}{#1}

\titleformat{\subsubsection}[block]
  {\sffamily\normalsize\bfseries\color{lpInk}}
  {}{0pt}{#1}

\titlespacing*{\section}{0pt}{14pt}{8pt}
\titlespacing*{\subsection}{0pt}{10pt}{4pt}
\titlespacing*{\subsubsection}{0pt}{8pt}{2pt}

% --- Hyperlinks ----------------------------------------------------------
\usepackage[hidelinks,unicode]{hyperref}

% --- TOC -----------------------------------------------------------------
\renewcommand{\contentsname}{\sffamily\Large\bfseries\color{lpAccent}Inhalt}

% --- Helfer --------------------------------------------------------------
\newcommand{\akzent}[1]{{\caveatfont\color{lpAmber}#1}}
\newcommand{\fachbegriff}[1]{\textbf{#1}}
\newcommand{\quelle}[1]{{\footnotesize\color{lpGray}(#1)}}

% =========================================================================
\begin{document}

% --- COVER ---------------------------------------------------------------
\begin{titlepage}
  \pagestyle{empty}
  \vspace*{1.5cm}
  \noindent{\sffamily\footnotesize\color{lpGray}LERNPLATT \textperiodcentered{} BY CAN SIEBERT}\\[2pt]
  \noindent{\color{lpAccent}\rule{\linewidth}{1.2pt}}\\[4pt]
  \noindent{\sffamily\footnotesize\color{lpGray}MODUL 4 \textperiodcentered{} TAG 8 \textperiodcentered{} KURS 22 (Bo 143) \textperiodcentered{} 9.\,Juni 2026}

  \vspace{3cm}
  {\sffamily\fontsize{30}{34}\selectfont\bfseries\color{lpInk}Dynamische\\Preisstrategien\\\& Performance\par}

  \vspace{0.7cm}
  {\sffamily\large\color{lpGray}Preis als Steuerungsinstrument --\\Markt, Marge, Marke \& Monitoring\\differenziert führen.\par}

  \vspace{1.5cm}
  {\caveatfont\LARGE\color{lpAmber}Lehrskript -- Tag 8 von 16}\par

  \vfill

  \noindent\begin{minipage}{0.55\linewidth}
    {\sffamily\footnotesize\color{lpGray}Lernpfad:}\\
    {\sffamily\small\color{lpInk}Wissen \textrightarrow{} Anwendung \textrightarrow{} Management}\\[10pt]
    {\sffamily\footnotesize\color{lpGray}Bearbeitungsdauer:}\\
    {\sffamily\small\color{lpInk}1 Kurstag (8 Unterrichtseinheiten)}
  \end{minipage}\hfill
  \begin{minipage}{0.4\linewidth}
    \raggedleft
    {\sffamily\footnotesize\color{lpGray}Dozent:}\\
    {\sffamily\small\color{lpInk}Can Siebert}\\[10pt]
    {\sffamily\footnotesize\color{lpGray}Format:}\\
    {\sffamily\small\color{lpInk}Theorie \textperiodcentered{} Fallstudie \textperiodcentered{} Coaching}
  \end{minipage}

  \vspace{0.5cm}
  \noindent{\color{lpAccent}\rule{\linewidth}{0.6pt}}
\end{titlepage}

% --- LERNZIELE -----------------------------------------------------------
\section*{Lernziele dieses Tages}
\addcontentsline{toc}{section}{Lernziele dieses Tages}
\thispagestyle{plain}

\noindent An Tag~7 haben Sie Sichtbarkeit, Content, Retention und Bewertungen als ein System verstanden. Heute ergänzt der Preis das Bild. \emph{Dynamische Preisstrategie} ist mehr als Rabattlogik -- sie ist ein zentrales Steuerungsinstrument für Marge, Conversion, Marktposition und Markenwirkung. Heute lernen Sie, Preise als Managementinstrument zu verwenden, ohne in einen Preiskrieg zu geraten.

\subsection*{L1 -- Wissen (Reproduktion und Einordnung)}
\begin{itemize}
  \item Grundlagen dynamischer Preisstrategien verstehen.
  \item Marktanalyse im Plattformverkauf einordnen.
  \item Einfluss von Wettbewerb, Nachfrage, Verfügbarkeit, Marge und Plattformranking kennen.
  \item Performance-Monitoring als Steuerungsinstrument verstehen.
\end{itemize}

\subsection*{L2 -- Anwendung (Transfer auf konkrete Situationen)}
\begin{itemize}
  \item Preis- und Marktinformationen auswerten.
  \item Dynamische Preisentscheidungen auf konkrete Plattformprodukte anwenden.
  \item Zielkonflikte zwischen Umsatz, Marge, Sichtbarkeit und Markenpositionierung bewerten.
  \item Plattformperformance anhand relevanter KPIs analysieren.
\end{itemize}

\subsection*{L3 -- Management (Entscheidung und Verantwortung)}
\begin{itemize}
  \item Preisstrategie als Teil einer integrierten Plattform-Verkaufsarchitektur gestalten.
  \item Entscheidungen unter Unsicherheit, Wettbewerbsdruck und Margendruck treffen.
  \item Verantwortung für Preis-Governance, Monitoring und strategische Steuerung übernehmen.
  \item Plattform-Verkauf als Zusammenspiel aus Markt, Preis, Content, Sichtbarkeit, Reputation und Profitabilität steuern.
\end{itemize}

\begin{merkbox}[Roter Faden]
Dynamische Preisstrategie ist \emph{kein} Synonym für Rabattaktion. Sie ist die Disziplin, Preise nach Produktrolle, Marktphase, Wettbewerb, Nachfrage und Marge \emph{differenziert} zu führen -- und sie an einer expliziten Governance auszurichten. Wer Preise pauschal senkt, weil Wettbewerber senken, hat keine Strategie. Er hat einen Reflex.
\end{merkbox}

\clearpage

% --- 1. EINSTIEG ---------------------------------------------------------
\section{Einstieg: Was dynamisch \emph{nicht} bedeutet}

In vielen Organisationen wird ``dynamische Preisgestaltung'' synonym mit ``Rabatte bei Wettbewerbsdruck'' verwendet. Das ist eine Verkürzung -- und eine teure. Dynamische Preisstrategie umfasst alle Mechanismen, in denen Preise sich an Wettbewerbssituation, Nachfrage, Lagerbestand, Saisonalität, Lieferfähigkeit und Marge orientieren. Rabatte sind \emph{ein} Werkzeug davon, kein Synonym. \quelle{Simon \& Fassnacht, 2015; Nagle \& Müller, 2016}

\subsection{Drei Stufen der Preissteuerung}

\begin{enumerate}
  \item \fachbegriff{Statische Preisgestaltung.} Preise werden einmal festgelegt und über lange Zeit gehalten. Robust, einfach -- und auf Marktplätzen meist nicht wettbewerbsfähig.
  \item \fachbegriff{Aktionsbasierte Preisgestaltung.} Preise werden in zeitlich begrenzten Aktionen verändert (Weihnachten, Black Friday, Sortimentswechsel). Planbar, geübt -- aber wenig flexibel auf laufende Marktveränderungen.
  \item \fachbegriff{Dynamische Preisstrategie.} Preise reagieren laufend (oder zumindest häufig) auf eine Mischung aus Wettbewerb, Nachfrage, Lager, Marge und Plattform-KPIs. Anspruchsvoll -- erfordert Governance, Daten und klare Grenzen.
\end{enumerate}

\subsection{Warum Plattformen dynamische Preise erzwingen}

Auf Plattformen sind Preise transparent vergleichbar. Wettbewerber ändern Preise teils mehrfach pro Woche. Algorithmen bevorzugen wettbewerbsfähige Angebote in Sichtbarkeit und Buy-Box. Wer mit statischen Preisen arbeitet, verliert -- nicht weil das Produkt schlecht wäre, sondern weil das System ihn unterhalb der Wahrnehmungsschwelle drückt. \quelle{Chen et al., 2016; Deshpandé \& Hwang, 2018}

\subsection{Warum dynamische Preise gefährlich werden können}

Dynamik ohne Governance erzeugt typische Schäden:

\begin{itemize}
  \item \textbf{Preiskrieg.} Eine Senkung führt zur Gegensenkung, das nächste Mal noch tiefer, bis die Marge weg ist.
  \item \textbf{Markenwahrnehmung.} Wer dauerhaft rabattiert, wird als ``Rabattmarke'' wahrgenommen -- nicht mehr als Qualitätsanbieter.
  \item \textbf{Kundenerwartung.} Kunden warten auf den nächsten Rabatt; Vollpreisverkäufe werden seltener.
  \item \textbf{Kannibalisierung.} Aktionen auf Standardprodukten ziehen Käufer von hochmargigen Premiumprodukten ab.
  \item \textbf{Vertrauensverlust.} Wechselnde Preise wirken bei B2B-Kunden unseriös.
\end{itemize}

\begin{eselsbox}[Eselsbrücke: \akzent{P-N-L-S-M} -- fünf Einflussfaktoren dynamischer Preise]
Fünf Buchstaben, fünf Einflüsse, die jede Preisentscheidung berücksichtigt:\\[2pt]
\textbf{P}reise der Wettbewerber\\
\textbf{N}achfrage \& Suchvolumen\\
\textbf{L}agerbestand \& Lieferfähigkeit\\
\textbf{S}aisonalität \& Plattformregeln\\
\textbf{M}arge \& Plattformgebühren\\[2pt]
\emph{Wer nur ``P'' beobachtet (Wettbewerber), reagiert. Wer alle fünf beobachtet, steuert.}
\end{eselsbox}

% --- 2. MARKTANALYSE IM PLATTFORMVERKAUF ---------------------------------
\section{Marktanalyse im Plattformverkauf}

Eine dynamische Preisstrategie ohne saubere Marktanalyse ist eine Bauchgefühlsentscheidung mit täglich neuer Begründung. Marktanalyse auf Plattformen umfasst sechs Dimensionen. \quelle{Simon \& Fassnacht, 2015; Porter, 1996}

\subsection{Sechs Dimensionen der Marktanalyse}

\begin{enumerate}
  \item \fachbegriff{Wettbewerberbeobachtung.} Welche Anbieter konkurrieren um die gleichen Keywords und Produktkategorien? Wie ist ihre Sortimentstiefe, ihre Preislage?
  \item \fachbegriff{Preisbandbreiten.} In welchem Korridor bewegen sich Wettbewerber? Wo liegt der ``Ankerpreis'' (Mittelwert)? Wo sind Premium-, wo Discount-Marken?
  \item \fachbegriff{Sortimentsvergleich.} Welche Funktionen, welcher Lieferumfang, welche Garantien stehen hinter den Preisen?
  \item \fachbegriff{Nachfrageentwicklung.} Wie verändern sich Suchvolumen, Klickrate, Conversion im Zeitverlauf?
  \item \fachbegriff{Bewertungslage.} Wie schneiden Wettbewerber bei Bewertungen ab -- besser, schlechter, gleich? Wo gibt es Bewertungslücken, die uns Spielraum geben?
  \item \fachbegriff{Sichtbarkeit.} Wer rangiert wo? Wer kauft sich nach oben, wer rankt organisch?
\end{enumerate}

\subsection{Marktanalyse als laufende Disziplin}

Marktanalyse ist kein Einmalprojekt. Sie ist eine \emph{laufende Disziplin} mit Reporting-Rhythmus (täglich/wöchentlich/monatlich, je nach Volatilität). Im einfachsten Fall ist sie eine Excel-Tabelle mit den fünf wichtigsten Wettbewerbern; im komplexeren Fall ein automatisiertes Tool, das Preisdifferenzen und Sichtbarkeitsverschiebungen meldet.

\subsection{Drei typische Anti-Muster}

\begin{itemize}
  \item \textbf{Nur den Marktführer beobachten.} Übersieht aggressive Newcomer und Nischenanbieter, die mit anderem Modell unterschneiden.
  \item \textbf{Preis nur isoliert vergleichen.} Ein 89-EUR-Wettbewerber ist nicht zwingend günstiger, wenn unser Lieferumfang doppelt so groß ist. Erst Vergleich auf Funktionsachse macht Preise vergleichbar.
  \item \textbf{Nachfrage ignorieren.} Wenn Suchvolumen für eine Kategorie sinkt, helfen Preissenkungen nicht -- der Markt wird kleiner, nicht der Anteil.
\end{itemize}

\begin{merkbox}[Marktanalyse-Faustregel]
Vergleichen Sie nie Preis ohne Funktion. Vergleichen Sie nie Funktion ohne Bewertung. Vergleichen Sie nie Bewertung ohne Nachfrageentwicklung. Erst die Kombination aller drei Achsen ergibt eine belastbare Marktsicht.
\end{merkbox}

% --- 3. PREISSTRATEGIEN -------------------------------------------------
\section{Preisstrategien auf Plattformen}

Ein dynamisches Preissystem nutzt nicht eine, sondern eine \emph{Mischung} aus Strategien -- typischerweise differenziert nach Produktgruppe.

\subsection{Acht Preisstrategien}

\begin{enumerate}
  \item \fachbegriff{Premiumpreis.} Bewusst über dem Marktdurchschnitt, gestützt durch Qualität, Service, Marke. Funktioniert nur, wenn der Mehrwert sichtbar ist.
  \item \fachbegriff{Wettbewerbsparität.} Preis am Marktdurchschnitt ausgerichtet. Sicher, aber wenig differenzierend.
  \item \fachbegriff{Penetrationspreis.} Bewusst unter dem Markt, um schnell Anteile zu gewinnen. Riskant -- erschwert späteres Anheben.
  \item \fachbegriff{Bündelpreis.} Hauptprodukt + Zubehör + Service als Paket; reduziert Vergleichbarkeit, erhöht Warenkorbwert.
  \item \fachbegriff{Mengenrabatt.} Skalierte Preise bei größeren Mengen; klassisch im B2B.
  \item \fachbegriff{Zeitlich begrenzte Aktion.} Bewusst befristete Senkung; aktiviert Aufmerksamkeit ohne dauerhaften Preisanker.
  \item \fachbegriff{Wertorientierter Preis.} Preis abgeleitet aus konkretem Kundennutzen (Einsparung, Umsatzwirkung). Im B2B besonders relevant.
  \item \fachbegriff{Differenzierter Preis nach Produktsegmenten.} Premium-Produkte stabil, Standard-Produkte dynamisch, Ersatzteile preisrobust -- siehe Sektion 4.
\end{enumerate}

\subsection{Welche Strategie wann?}

\begin{tabularx}{\linewidth}{@{}lXXX@{}}
\toprule
\textbf{Strategie} & \textbf{Wann passend} & \textbf{Risiko} & \textbf{KPI-Fokus} \\
\midrule
Premiumpreis & starke Marke, hoher Mehrwert & Sichtbarkeitsverlust & Marge, Bewertungsdurchschnitt \\
Wettbewerbsparität & generische Produkte, Konsolidierung & Verlust Profil & Umsatz, Marktanteil \\
Penetration & Markteintritt, neue Plattform & Margenverbrennung & Marktanteil, CAC \\
Bündelpreis & Vergleichsdruck hoch, Cross-Sell möglich & Bundle-Komplexität & Warenkorbwert, Deckungsbeitrag \\
Wertorientiert & B2B, klar quantifizierbarer Nutzen & Argumentationsaufwand & Marge, Lifetime Value \\
\bottomrule
\end{tabularx}

\subsection{Risiken dynamischer Preisgestaltung im Detail}

\begin{itemize}
  \item \textbf{Margenverlust.} Jede Preissenkung kostet -- pro Verkauf, multipliziert mit dem Volumen.
  \item \textbf{Preiskrieg.} Spiralartige gegenseitige Senkungen, die niemand gewinnt.
  \item \textbf{Markenwahrnehmung.} Permanent niedrige Preise senken die wahrgenommene Qualität.
  \item \textbf{Kundenerwartung.} ``Bei Ihnen ist immer Rabatt'' -- Vollpreisverkäufe werden seltener.
  \item \textbf{Kannibalisierung.} Aktionen ziehen Käufer von Premiumprodukten ab.
  \item \textbf{Vertrauensverlust.} Wechselnde Preise wirken unseriös, besonders im B2B.
\end{itemize}

\quelle{vgl.\ Simon \& Fassnacht, 2015; Shapiro \& Varian, 1999}

% --- 4. DIFFERENZIERUNG NACH PRODUKTROLLE --------------------------------
\section{Differenzierung nach Produktrolle}

Eine zentrale Senior-Erkenntnis: \emph{Nicht alle Produkte sollten gleich bepreist werden}. Eine reife Plattform-Verkaufsorganisation unterscheidet typischerweise vier Produktrollen mit jeweils eigener Preislogik.

\subsection{Vier Produktrollen}

\begin{enumerate}
  \item \fachbegriff{Premium-/Markenprodukte.} Hohe Marge, strategische Bedeutung, Markenwirkung. \emph{Preislogik:} stabil halten, über Wert kommunizieren, keine pauschalen Rabatte. Bei Wettbewerbsdruck antwortet Content und Service -- nicht der Preis.
  \item \fachbegriff{Standardprodukte mit Wettbewerbsdruck.} Hohe Vergleichbarkeit, hohe Preiselastizität. \emph{Preislogik:} gezielte zeitlich begrenzte Aktionen mit Mindestmarge; gegebenenfalls Wettbewerbsparität.
  \item \fachbegriff{Ersatzteile / wiederkehrende Käufe.} Geringe Vergleichbarkeit (oft Kompatibilitäts-gebunden), hohe Wiederkaufrate. \emph{Preislogik:} stabile Preise; Differenzierung über Verfügbarkeit, schnelle Lieferung, Kompatibilitätsinformation.
  \item \fachbegriff{Bundles.} Aus zwei oder mehr Komponenten kombiniert. \emph{Preislogik:} bewusster Bundle-Vorteil (z.\,B.\ 10--15\,\% gegenüber Einzelpreissumme); reduziert Vergleichbarkeit und erhöht Warenkorbwert.
\end{enumerate}

\subsection{Beispiel: differenzierte Preisstrategie}

Für ``KitchenPro Direct'' (siehe Sektion 9) sieht eine differenzierte Strategie typischerweise so aus:

\begin{itemize}
  \item Premium-Küchenmaschinen \textrightarrow{} \emph{preisstabil}, mit Service-Bundles und Anwendungs-Content angereichert.
  \item Standardzubehör \textrightarrow{} \emph{dynamisch} mit zeitlich begrenzten Aktionen, Mindestmarge geschützt.
  \item Ersatzteile \textrightarrow{} \emph{stabil}, Differenzierung über Verfügbarkeit und schnelle Lieferung.
  \item Bundles \textrightarrow{} \emph{aktiv ausgebaut}, um Vergleichbarkeit zu reduzieren.
\end{itemize}

\begin{merkbox}[Differenzierungs-Prinzip]
Eine reife Preisstrategie behandelt verschiedene Produktrollen unterschiedlich. Wer alle Produkte gleich rabattiert, beschädigt Premium-Marke und Margen-Architektur gleichzeitig. Wer keine Produktrolle definiert, hat keine differenzierte Preisstrategie -- sondern einen Preisreflex.
\end{merkbox}

% --- 5. PREIS-GOVERNANCE ------------------------------------------------
\section{Preis-Governance}

Dynamische Preise ohne Governance führen zu Margenverlusten und Markenschäden. Daher: jedes dynamische Preissystem braucht klare Regeln, Grenzen und Entscheidungskompetenzen.

\subsection{Sechs Governance-Bausteine}

\begin{enumerate}
  \item \fachbegriff{Mindestmarge.} Pro Produktgruppe ein Mindestwert (z.\,B.\ 25\,\%), unter den nicht gegangen wird -- auch nicht ``für nur eine Woche''.
  \item \fachbegriff{Preiskorridor.} Maximaler Abstand zum Marktdurchschnitt nach oben und unten.
  \item \fachbegriff{Aktionsbudget.} Wie viel Marge geben wir pro Quartal für Aktionen frei?
  \item \fachbegriff{Entscheidungsrechte.} Wer darf welche Preisveränderung freigeben? Algorithmus, Pricing Manager, Geschäftsführung?
  \item \fachbegriff{Eskalationsregeln.} Wann muss eine Preisentscheidung in die Geschäftsführung -- z.\,B.\ Dauerunterschreitung, Preiskampf, Markenrisiko?
  \item \fachbegriff{Lernschleifen.} Quartalsweise Auswertung: Welche Aktion hat Marge gebracht, welche hat sie nur verschoben?
\end{enumerate}

\subsection{Wann manuelle Governance algorithmische Preise überschreibt}

Algorithmen können viel -- aber sie sollten in bestimmten Situationen \emph{nicht} alleine entscheiden:

\begin{itemize}
  \item Beim Auftreten neuer Wettbewerber mit unklarer Strategie (Reaktion abwarten, Daten sammeln).
  \item Bei Markenwirkungsrisiken (z.\,B.\ Premium-Produkt droht in Discount-Wahrnehmung zu rutschen).
  \item Bei Vorzeichen eines Preiskriegs (Spirale erkennen und unterbrechen).
  \item Bei extremer Nachfrageschwankung (Algorithmus überreagiert auf temporäre Effekte).
\end{itemize}

\subsection{Algorithmisches Pricing -- Chancen und Grenzen}

Algorithmische Preissysteme können sekundenschnell reagieren -- z.\,B.\ auf Amazon-Marktplätzen mehrfach pro Stunde. Empirische Studien zeigen aber, dass diese Systeme \emph{Preiskriege} verstärken können, wenn nicht gegengesteuert wird. \quelle{Chen et al., 2016; Deshpandé \& Hwang, 2018}

\begin{eselsbox}[Eselsbrücke: \akzent{M-K-A-E-E-L} -- sechs Bausteine der Preis-Governance]
Sechs Buchstaben, sechs Bausteine:\\[2pt]
\textbf{M}indestmarge -- nicht verhandelbar.\\
\textbf{K}orridor -- definierter Bereich.\\
\textbf{A}ktionsbudget -- begrenzte ``Margenfreigabe''.\\
\textbf{E}ntscheidungsrechte -- klar zugeordnet.\\
\textbf{E}skalation -- wann geht es ins Management?\\
\textbf{L}ernschleifen -- quartalsweise auswerten.\\[2pt]
\emph{Wer auch nur einen der sechs nicht festlegt, hat keine Preisstrategie -- nur eine Preistabelle, die manchmal jemand ändert.}
\end{eselsbox}

% --- 6. PERFORMANCE-MONITORING -------------------------------------------
\section{Performance-Monitoring}

Eine integrierte Plattform-Verkaufsstrategie braucht Steuerung -- und Steuerung braucht KPIs, die nicht nur Aktivität, sondern Wirkung messen. \quelle{Kaplan \& Norton, 1996; Ellis \& Brown, 2017}

\subsection{Operative vs.\ strategische KPIs}

\begin{itemize}
  \item \textbf{Operative KPIs.} Tagesgeschäft -- Klicks, Klickrate, Aufrufe, Sichtbarkeit-Rank, Lagerstand.
  \item \textbf{Taktische KPIs.} Monat -- Conversion Rate, Warenkorbwert, Retourenquote, Bewertungsentwicklung, Aktionsrentabilität.
  \item \textbf{Strategische KPIs.} Quartal / Jahr -- Deckungsbeitrag pro Produktgruppe, Marge pro Kanal, Wiederkaufrate, Plattformabhängigkeit.
\end{itemize}

\subsection{Kern-KPIs für Plattform-Performance}

\begin{enumerate}
  \item \fachbegriff{Umsatz pro Produktgruppe.}
  \item \fachbegriff{Deckungsbeitrag pro Produktgruppe.}
  \item \fachbegriff{Conversion Rate je Produktgruppe.}
  \item \fachbegriff{Preisabstand zu Hauptwettbewerbern.}
  \item \fachbegriff{Klickrate und Sichtbarkeit-Rank.}
  \item \fachbegriff{Warenkorbwert.}
  \item \fachbegriff{Wiederkaufrate.}
  \item \fachbegriff{Retourenquote.}
  \item \fachbegriff{Bewertungsentwicklung.}
  \item \fachbegriff{Aktionsrentabilität} -- gemessen \emph{nach Marge}, nicht nur nach Umsatz.
\end{enumerate}

\subsection{Steuerungslogik: messen, bewerten, entscheiden, testen}

Eine reife Performance-Steuerung folgt vier Schritten:

\begin{enumerate}
  \item \textbf{Messen.} KPIs in definiertem Rhythmus erfassen.
  \item \textbf{Bewerten.} Abweichungen gegen Ziel oder Vorperiode prüfen -- mit Kontext.
  \item \textbf{Entscheiden.} Bei kritischer Abweichung Aktion einleiten; bei stabilem Ergebnis: keine Aktion ist auch eine Entscheidung.
  \item \textbf{Testen.} Hypothesen über Experimente prüfen (z.\,B.\ Preisexperimente bei zwei Standardprodukten über 4 Wochen).
\end{enumerate}

\subsection{Falle: Vanity-Metriken}

\fachbegriff{Vanity-Metriken} sind Größen, die imposant aussehen, aber nichts steuern: Impressionen, Followerzahlen, ``Reichweite''. Senior-Level-Reports sollten sie weglassen oder nur als Nebenangabe behandeln -- die Steuerung erfolgt über Marge, Deckungsbeitrag und Wiederkaufrate.

\begin{merkbox}[Monitoring-Faustregel]
Wenn eine KPI im Managementreport keine Entscheidung auslöst, gehört sie nicht in den Report. Sinn von Monitoring ist nicht, Aktivität nachzuweisen -- sondern Entscheidungen zu ermöglichen.
\end{merkbox}

% --- 7. INTEGRATION: SEO x CONTENT x PREIS x BEWERTUNG -------------------
\section{Integration: Plattform-Verkauf als ein System}

Tag~7 hat Sichtbarkeit, Content, Retention und Bewertungen behandelt. Tag~8 ergänzt Markt und Preis. Erst zusammen entsteht eine integrierte Plattform-Verkaufsarchitektur. \quelle{Kotler et al., 2017; Homburg, 2020}

\subsection{Sechs Hebel im Verbund}

\begin{enumerate}
  \item \textbf{SEO} -- Sichtbarkeit.
  \item \textbf{Content} -- Conversion-Vorbereitung.
  \item \textbf{Preis} -- Conversion-Hebel und Marketing-Signal zugleich.
  \item \textbf{Bewertungen} -- Vertrauensbeschleuniger.
  \item \textbf{Retention} -- Marge über Wiederkauf.
  \item \textbf{Monitoring} -- Steuerungsgrundlage.
\end{enumerate}

\subsection{Wenn eine Stellschraube wackelt}

Jede der sechs Stellschrauben kann eine andere kippen lassen:

\begin{itemize}
  \item Aggressive Preissenkung \textrightarrow{} mehr Klicks, aber falsche Kundenerwartung \textrightarrow{} schlechte Bewertungen \textrightarrow{} Ranking-Verlust.
  \item Schlechter Content \textrightarrow{} niedrige Conversion \textrightarrow{} Algorithmus rankt herunter \textrightarrow{} Sichtbarkeit weg.
  \item Schwacher Bewertungen-Prozess \textrightarrow{} Sterne sinken \textrightarrow{} Conversion sinkt \textrightarrow{} höherer CAC -- Marge weg.
\end{itemize}

\subsection{Was eine integrierte Architektur leistet}

\begin{itemize}
  \item Pro Produktgruppe definierte Preisrolle (siehe Sektion 4).
  \item Pro Produktgruppe definierte Keyword- und Content-Standards (Tag 7).
  \item Bewertungsmanagement, das in beide Richtungen wirkt: für Sichtbarkeit \emph{und} als Frühwarnsystem.
  \item Retention-Mechanismus, der Bestandskunden hält und Marge stabilisiert.
  \item Monitoring-Rhythmus, der Preis-, Conversion-, Bewertungs- und Margenentwicklung verbindet.
\end{itemize}

% --- 8. FALLSTUDIE-LEITFADEN ---------------------------------------------
\section{Fallstudie-Leitfaden: KitchenPro Direct}

In den Unterrichtseinheiten 4--5 bearbeiten Sie die Fallstudie ``KitchenPro Direct'' -- ein hochwertiger Anbieter für professionelle Gastronomie und ambitionierte Privatkunden, mit 6{,}5\,Mio.\,EUR Plattformumsatz, 29\,\% Marge, 4{,}6 Sternen und einer sinkenden Conversion Rate.

\subsection{Analytischer Aufbau}

\begin{enumerate}
  \item \textbf{Markt- und Wettbewerbssituation.} Was steht im Markt -- wer rabattiert, wer nicht? Wo sind Sortimentslücken?
  \item \textbf{Produktrollen bewerten.} Premium-Küchenmaschinen, Standardzubehör, Ersatzteile, Bundles -- mit welcher Preisrolle?
  \item \textbf{Differenzierte Preisstrategie.} Pro Produktgruppe: stabil, dynamisch, gebündelt oder premium-positioniert?
  \item \textbf{Performance-Monitoring.} KPIs pro Produktgruppe, mit klarer Schwelle für Eingriff.
  \item \textbf{Budgetverteilung.} Wie verteilen sich 120.000\,EUR auf Monitoring, Content, Bundles, Tests und Reporting?
  \item \textbf{Entscheidung unter Unsicherheit.} Welche kurzfristig attraktive Preisaktion wird trotz Wettbewerbsdruck \emph{nicht} gefahren?
\end{enumerate}

\subsection{Erwartete Empfehlungsrichtung}

In einer typischen Musterlösung wird differenziert: Premium-Geräte preisstabil mit Service-Bundles und Content-Stärkung; Standardzubehör dynamisch mit Mindestmarge; Ersatzteile stabil mit Verfügbarkeits-Differenzierung; Bundles aktiv ausgebaut. Die Entscheidung unter Unsicherheit: ``Trotz sinkender Conversion verzichten wir auf pauschale Rabatte bei Premium-Küchenmaschinen, weil Marke und Marge strategische Vermögenswerte sind -- stattdessen testen wir Bundles und investieren in Content.''

\begin{merkbox}[Argumentationsmuster für die Fallstudie]
Eine gute Preisstrategie auf einer Plattform enthält:
\begin{enumerate}
  \item \textbf{Produktrollen.} Pro Gruppe definiert.
  \item \textbf{Preislogik je Rolle.} Konsistent und begründet.
  \item \textbf{Governance.} Mindestmarge, Korridor, Entscheidungsrechte.
  \item \textbf{Monitoring.} KPIs mit Eingriffsschwelle.
  \item \textbf{Verzichtsentscheidung.} Mindestens eine bewusste Senkung, die wir \emph{nicht} machen.
\end{enumerate}
\end{merkbox}

% --- 9. MANAGEMENT-PERSPEKTIVE ------------------------------------------
\section{Management-Perspektive: Zielkonflikte}

Auf Wissens- und Anwendungsebene sieht eine gute Preisstrategie wie eine sauber differenzierte Tabelle aus. Auf Managementebene geht es um Zielkonflikte, die nicht ``gelöst'', sondern \emph{entschieden} werden müssen. \quelle{vgl.\ Porter, 1996; Simon \& Fassnacht, 2015}

\subsection{Fünf zentrale Zielkonflikte}

\begin{enumerate}
  \item \textbf{Umsatz vs.\ Marge.} Aggressive Preissenkung erhöht Umsatz, senkt Marge. Wann lohnt sich der Wechsel, wann nicht?
  \item \textbf{Sichtbarkeit vs.\ Markenpositionierung.} Plattform-Algorithmen belohnen niedrige Preise und hohe Conversion -- aber wer Premium ist, kann hier nicht beliebig folgen.
  \item \textbf{Kurzfristige Conversion vs.\ Preisanker.} Eine Aktion erhöht Conversion sofort -- aber sie verändert den Preisanker im Kundengedächtnis dauerhaft.
  \item \textbf{Wettbewerberreaktion vs.\ eigene Strategie.} Reflexartiges Mitziehen mit Wettbewerbern führt in einen Preiskrieg. Nicht-Reagieren kann aber kurzfristig Marktanteil kosten.
  \item \textbf{Algorithmus vs.\ Governance.} Algorithmen können schnell -- aber sollten in markenrelevanten Situationen nicht allein entscheiden.
\end{enumerate}

\subsection{Senior-Level-Test}

Eine senior-level-fähige Preisstrategie lässt sich in einem Satz erklären, der eine Verzichtsentscheidung enthält. Beispiel: ``Wir halten Premium-Preise stabil und stärken Content und Bundles -- wir verzichten bewusst auf pauschale Rabatte trotz sinkender Conversion, weil eine pauschale Senkung Marke und Margenarchitektur dauerhaft beschädigen würde.''

\subsection{Preisstrategie als Senior-Level-Verantwortung}

Preisentscheidungen wirken wie operativ -- sie sind aber strategisch. Sie verändern Kundenerwartung, Marke, Marge und Wettbewerbsverhalten auf Jahre. Wer Preis an die Pricing-Abteilung allein delegiert, delegiert Markenführung und Margenarchitektur. Beides gehört auf die Senior-Ebene.

\begin{eselsbox}[Eselsbrücke: \akzent{U-S-K-W-A}]
Fünf Zielkonflikte, die im dynamischen Pricing auf Senior-Level entschieden werden:\\[2pt]
\textbf{U}msatz vs.\ Marge.\\
\textbf{S}ichtbarkeit vs.\ Markenpositionierung.\\
\textbf{K}urzfrist (Conversion) vs.\ Preisanker.\\
\textbf{W}ettbewerberreaktion vs.\ eigene Strategie.\\
\textbf{A}lgorithmus vs.\ Governance.\\[2pt]
\emph{Fünf Konflikte, die kein Algorithmus allein verantworten sollte.}
\end{eselsbox}

\begin{merkbox}[Schluss-Merksatz für Tag 8]
Dynamische Preisstrategie ist kein automatisiertes Rabattsystem. Sie ist die Disziplin, Preise differenziert nach Produktrolle, Marktphase und Marge zu führen -- mit klarer Governance, sauberem Monitoring und bewussten Verzichtsentscheidungen. \emph{Wer Wettbewerberpreisen reflexartig folgt, verliert Marge, Marke und mittelfristig Position.}
\end{merkbox}

% --- 10. LITERATUR -------------------------------------------------------
\clearpage
\section{Literatur}

Im Skript verwendete und für die Vertiefung empfohlene Quellen. Zitierstil: APA-7.

\begin{description}[style=nextline,leftmargin=18pt,labelindent=0pt,itemsep=4pt]
  \item[Chen, L., Mislove, A., \& Wilson, C. (2016).] An empirical analysis of algorithmic pricing on Amazon Marketplace. In \emph{Proceedings of the 25th International Conference on World Wide Web} (S.\ 1339--1349). \href{https://doi.org/10.1145/2872427.2883089}{https://doi.org/10.1145/2872427.2883089}
  \item[Deshpandé, R., \& Hwang, J. (2018).] Dynamic pricing in digital markets. \emph{MIT Sloan Management Review, 60}(1), 1--9.
  \item[Ellis, S., \& Brown, M. (2017).] \emph{Hacking growth: How today's fastest-growing companies drive breakout success}. Crown Business.
  \item[Kaplan, R.\,S., \& Norton, D.\,P. (1996).] \emph{The balanced scorecard: Translating strategy into action}. Harvard Business School Press.
  \item[Nagle, T.\,T., \& Müller, G. (2016).] \emph{The strategy and tactics of pricing: A guide to growing more profitably} (6.\ Aufl.). Routledge.
  \item[Porter, M.\,E. (1996).] What is strategy? \emph{Harvard Business Review, 74}(6), 61--78.
  \item[Shapiro, C., \& Varian, H.\,R. (1999).] \emph{Information rules: A strategic guide to the network economy}. Harvard Business School Press.
  \item[Simon, H., \& Fassnacht, M. (2015).] \emph{Preismanagement: Strategie -- Analyse -- Entscheidung -- Umsetzung} (4.\ Aufl.). Springer Gabler.
\end{description}

% --- 11. GLOSSAR ---------------------------------------------------------
\clearpage
\section{Glossar}

Die wichtigsten Begriffe des Tages -- kurz definiert, in alphabetischer Reihenfolge.

\begin{description}[style=nextline,leftmargin=0pt,labelindent=0pt,itemsep=4pt]
  \item[Aktionsrentabilität] Bewertung einer Preisaktion nach Margenwirkung -- nicht nur nach Umsatzsteigerung.
  \item[Algorithmisches Pricing] Software-gestützte Preissetzung, die auf Wettbewerb, Nachfrage und Lager reagiert; kann ohne Governance Preiskriege verstärken.
  \item[Buy-Box] Bevorzugte Verkaufsbox auf einer Plattform; Sichtbarkeitsvorteil, beeinflusst durch Preis, Lieferfähigkeit, Verkäuferperformance.
  \item[Bundle (Preisbundle)] Kombination aus Hauptprodukt, Zubehör und/oder Service zu einem Bündelpreis.
  \item[Customer Acquisition Cost (CAC)] Gesamtkosten zur Gewinnung eines neuen Kunden.
  \item[Deckungsbeitrag] Differenz aus Umsatz und variablen Kosten; zentrale Profitabilitätsgröße pro Produkt oder Gruppe.
  \item[Differenzierte Preisstrategie] Unterschiedliche Preislogik je Produktrolle (Premium, Standard, Ersatz, Bundle).
  \item[Dynamische Preisstrategie] Laufende oder häufige Anpassung von Preisen an Wettbewerb, Nachfrage, Marge, Lager und Plattform-KPIs.
  \item[Eskalationsregel] Vorschrift, wann eine Preisentscheidung in die nächsthöhere Entscheidungsebene gehoben werden muss.
  \item[Governance (Preis)] Regelwerk für Preisentscheidungen: Mindestmarge, Korridor, Entscheidungsrechte, Aktionsbudget, Eskalation.
  \item[Kannibalisierung] Effekt, bei dem Aktionen auf einer Produktgruppe Käufer aus einer anderen Produktgruppe abziehen.
  \item[Marktanalyse (Plattform)] Strukturierte, laufende Beobachtung von Wettbewerb, Preisbandbreiten, Sortiment, Nachfrage, Bewertung und Sichtbarkeit.
  \item[Mindestmarge] Pro Produktgruppe definierter unterer Margenwert, der nicht unterschritten wird.
  \item[Penetrationspreis] Bewusst unter dem Markt liegende Einführung, um schnell Anteile zu gewinnen.
  \item[Performance-Monitoring] Strukturierte Erfassung und Auswertung relevanter KPIs zur Steuerung des Plattform-Verkaufs.
  \item[Preisanker] Vom Kunden wahrgenommener Referenzpreis, der durch frühere Preise und Wettbewerb geprägt wird.
  \item[Preisbandbreite] Korridor, in dem sich die Wettbewerberpreise innerhalb einer Produktkategorie bewegen.
  \item[Preiselastizität] Veränderung der Nachfrage in Reaktion auf eine Preisänderung; je höher, desto preissensibler die Kunden.
  \item[Preiskorridor] Definierter Spielraum, in dem dynamische Anpassungen erlaubt sind.
  \item[Preiskrieg] Spirale gegenseitiger Preissenkungen, die alle Beteiligten Marge kostet.
  \item[Premiumpreis] Preis bewusst über dem Marktdurchschnitt, gestützt durch Qualität, Service, Marke.
  \item[Produktrolle] Strategische Funktion eines Produkts im Sortiment (Premium, Standard, Ersatz, Bundle); bestimmt die Preislogik.
  \item[Vanity-Metrik] Kennzahl, die imposant wirkt, aber keine Steuerungsentscheidung auslöst (z.\,B.\ reine Impressionen).
  \item[Wertorientierter Preis] Preisbildung aus dem konkret quantifizierbaren Kundennutzen (Einsparung, Umsatzwirkung).
  \item[Wettbewerbsparität] Preis am Marktdurchschnitt orientiert; sicher, aber wenig differenzierend.
  \item[Zielkonflikt (Preis)] Spannungsfeld zwischen Umsatz, Marge, Sichtbarkeit, Markenpositionierung und Wettbewerberreaktion -- nicht ``lösbar'', sondern bewusst entscheidbar.
\end{description}

\vspace{1cm}
\noindent{\color{lpAccent}\rule{\linewidth}{0.6pt}}\\[4pt]
{\footnotesize\sffamily\color{lpGray}Lernplatt \textperiodcentered{} by Can Siebert \textperiodcentered{} Kurs 22 (Bo 143) \textperiodcentered{} Tag 8 \textperiodcentered{} \textcopyright{} 2026}

\end{document}
